\label{sec:set-operations-on-real-line}

\begin{topicbox}{Set Operations as Real-Line Bookkeeping}
\begin{exposition}
The set operations from Volume I are now used with subsets of \(\mathbb R\).
This section does not introduce new set theory. It fixes the notation and the
finite-versus-arbitrary distinction that later controls neighborhoods, open
sets, closed sets, covers, and compactness.
\end{exposition}
\end{topicbox}

\begin{remark*}[Notation]
An indexed family
\[
  (E_i)_{i\in I}
\]
is read here as a family of real-line subsets when
\[
  \forall i\in I\;(E_i\subseteq\mathbb R).
\]
The ambient set is fixed to be \(\mathbb R\). Unions, intersections,
complements, and differences are therefore read as operations on subsets of
the real line, not as new arithmetic operations on real numbers.
\end{remark*}
\begin{dependencies}
\begin{itemize}
  \item \hyperref[def:indexed-family]{Indexed Family of Sets}
  \item \hyperref[def:subset]{Subset}
\end{itemize}
\end{dependencies}

\begin{remark*}[Unions and intersections]
For a family \((E_i)_{i\in I}\) of real-line subsets,
\[
  \bigcup_{i\in I}E_i
  =
  \{x\in\mathbb R:\exists i\in I\ (x\in E_i)\},
\]
and
\[
  \bigcap_{i\in I}E_i
  =
  \{x\in\mathbb R:\forall i\in I\ (x\in E_i)\}.
\]
Thus a union records membership in at least one set, while an intersection
records simultaneous membership in every set.
\end{remark*}

\begin{figure}[htbp]
\centering
\begin{tikzpicture}[atlas, x=1cm, y=1cm]
  \draw[atlas axes] (-0.25,0) -- (7.1,0);
  \draw[atlas axes] (-0.25,-1.35) -- (7.1,-1.35);
  \draw[atlas axes] (-0.25,-2.7) -- (7.1,-2.7);

  \draw[atlas curve, atlasblue, line width=2pt] (0.7,0) -- (2.7,0);
  \draw[atlas curve, atlasgreen, line width=2pt] (1.9,0.16) -- (4.2,0.16);
  \draw[atlas curve, atlasorange, line width=2pt] (3.7,-0.16) -- (6.1,-0.16);
  \node[atlas label, left] at (-0.32,0) {$E_i$};

  \draw[atlas curve, atlasblue, line width=2.4pt] (0.7,-1.35) -- (6.1,-1.35);
  \node[atlas label, left] at (-0.32,-1.35) {$\bigcup_i E_i$};

  \draw[atlas curve, atlasmagenta, line width=2.4pt] (3.7,-2.7) -- (4.2,-2.7);
  \node[atlas label, left] at (-0.32,-2.7) {$\bigcap_i E_i$};

  \node[atlas label, below] at (0.7,-0.18) {$a$};
  \node[atlas label, below] at (6.1,-0.18) {$b$};
  \node[atlas label, right] at (6.4,-1.35) {covered by at least one set};
  \node[atlas label, right] at (4.45,-2.7) {common part};
\end{tikzpicture}

\caption{A family of real-line subsets, its union, and its common intersection.}
\label{fig:set-operations-on-real-line}
\end{figure}

\begin{remark*}[Finite versus arbitrary operations]
The distinction between finite and arbitrary operations is structural. Later,
open sets will be stable under arbitrary unions but only finite intersections.
Closed sets will satisfy the dual pattern: arbitrary intersections but only
finite unions. The words finite and arbitrary therefore mark which operations
survive the topological definitions.
\end{remark*}

\begin{remark*}[Complements and De Morgan laws]
Complements are always taken inside \(\mathbb R\):
\[
  E^c=\mathbb R\setminus E.
\]
For a real-line family \((E_i)_{i\in I}\), the indexed De Morgan laws become
\[
  \left(\bigcup_{i\in I}E_i\right)^c
  =
  \bigcap_{i\in I}E_i^c,
  \qquad
  \left(\bigcap_{i\in I}E_i\right)^c
  =
  \bigcup_{i\in I}E_i^c.
\]
These are the set-theoretic identities that later convert open-set closure
laws into closed-set closure laws.
\end{remark*}
\begin{dependencies}
\begin{itemize}
  \item \hyperref[def:union]{Union}
  \item \hyperref[def:intersection]{Intersection}
  \item \hyperref[def:complement]{Complement}
  \item \hyperref[def:indexed-union]{Indexed Union}
  \item \hyperref[def:indexed-intersection]{Indexed Intersection}
  \item \hyperref[thm:indexed-de-morgan]{De Morgan's Laws for Indexed Families}
\end{itemize}
\end{dependencies}

\begin{remark*}[Forward use]
Neighborhoods and open sets use unions to assemble local pieces into larger
sets. Covers use families whose union contains the set being studied.
Compactness asks whether an arbitrary covering family can be replaced by a
finite subfamily. The same set operations recur throughout the rest of the
chapter.
\end{remark*}
